\documentclass[12pt, twoside]{report}

% Paquetes LaTeX y estilos globales
\input{configuracion/paquetes}
\input{configuracion/estilo} 

% Con \hypersetup podremos introducir los metadatos para la descripción del documento PDF. Recuerda modificarlo cuando transformes el documento.
\hypersetup{ 
    pdfauthor={Autor del trabajo}
    pdftitle={Título del trabajo}
    pdfsubject={Plantilla de la Universidad de Cádiz}
    pdfkeywords{Palabras clave del documento}
    pdflang={Idioma del documento}
    pdfcreator={Creador del documento}
}

% Comienzo del documento
\begin{document}

    % Portada y secciones no numeradas
    \input{sections/00_portada}
    \input{sections/00_agradecimientos}
    \input{sections/00_declaracion}
    \input{sections/00_resumen}
    \input{sections/00_abstract}
    
    % Índice del documento y de figuras
    \begingroup
        % Los enlaces son normalmente azules, pero en los índices se configuran a negro
        % para que no aparezca todo azul
        \hypersetup{linkcolor=black}
\newpage{\ }
\thispagestyle{empty}
\tableofcontents
        \listoffigures
        \listoftables
    \endgroup
    
    % Cambia el estilo de números de página a normal
    \clearpage\pagenumbering{arabic}
    
    % Capítulos del trabajo
    % IMPORTANTE: POR % delante de la siguiente linea para que no aparezca en el trabajo final
    \input{sections/Ejemplos_comandos_latex}
    \input{sections/01_Introduccion}
    \input{sections/02_Marco Teórico}
    \input{sections/03_Metodología}
    \input{sections/04_Resultados}
    \input{sections/05_Discusión}
    \input{sections/XX_Conclusiones}

    % Bibliografía
    \newpage{\ }
    \thispagestyle{empty}
    \setcitestyle{aysep={}}
    \bibliographystyle{apalike7}
    \bibliography{bibliografia}

    %Anexos
    \input{sections/YY_Anexos}

    

% Fin del documento
\end{document}
